\documentclass[aspectratio=169]{beamer}
\usepackage[utf8]{inputenc}
\usepackage{lmodern}
\usepackage[T1]{fontenc}
\usepackage{listings}
\usepackage{tikz}
\usetikzlibrary{arrows.meta, positioning, shapes.geometric, calc, graphs}

% Set path to find the theme
\makeatletter
\def\input@path{{../BeamerTemplate/theme/}}
\makeatother

\usetheme{CleanEasy}
\input{../BeamerTemplate/configs/configs}

% Listings configuration
\lstset{
    basicstyle=\ttfamily\footnotesize,
    keywordstyle=\color{blue},
    commentstyle=\color{gray},
    stringstyle=\color{red},
    showstringspaces=false,
    breaklines=true,
    frame=single,
    numbers=left,
    numberstyle=\tiny\color{gray}
}

\title{Graph Algorithms}
\subtitle{Compute Connectivity, Paths, and Structures Over Graphs}
\author{Minseok Jeon}
\institute{DGIST}
\date{\today}

\begin{document}

\begin{frame}
  \titlepage
\end{frame}

\begin{frame}{Table of Contents}
  \tableofcontents
\end{frame}

\section{Introduction}

\begin{frame}{What are Graph Algorithms?}
  \textbf{Graph Algorithms}: Methods to solve problems on graph structures

  \vspace{0.3cm}
  \textbf{Key Problems:}
  \begin{itemize}
    \item \textbf{Traversal}: Visit all vertices/edges
    \item \textbf{Connectivity}: Determine if vertices are connected
    \item \textbf{Shortest Paths}: Find minimum cost paths
    \item \textbf{Spanning Trees}: Connect all vertices with minimum cost
    \item \textbf{Flow}: Maximize throughput in networks
  \end{itemize}

  \vspace{0.3cm}
  \textbf{Real-World Applications:}
  \begin{itemize}
    \item Social networks (connections, recommendations)
    \item Navigation systems (GPS, routing)
    \item Computer networks (packet routing, topology)
    \item Task scheduling (dependencies, ordering)
    \item Bioinformatics (protein interactions, gene networks)
  \end{itemize}
\end{frame}

\begin{frame}{Graph Basics Recap}
  \textbf{Graph}: $G = (V, E)$ where $V$ = vertices, $E$ = edges

  \vspace{0.3cm}
  \textbf{Types:}
  \begin{itemize}
    \item \textbf{Directed vs Undirected}: Edges have direction or not
    \item \textbf{Weighted vs Unweighted}: Edges have costs or not
    \item \textbf{Cyclic vs Acyclic}: Contains cycles or not (DAG)
    \item \textbf{Connected vs Disconnected}: All vertices reachable or not
  \end{itemize}

  \vspace{0.3cm}
  \textbf{Complexity Notation:}
  \begin{itemize}
    \item $V$ = number of vertices
    \item $E$ = number of edges
    \item Dense graph: $E \approx V^2$
    \item Sparse graph: $E \approx V$
  \end{itemize}
\end{frame}

\section{BFS/DFS Applications}

\begin{frame}{Breadth-First Search (BFS)}
  \textbf{Level-by-Level Exploration Using Queue}

  \vspace{0.3cm}
  \textbf{Algorithm:}
  \begin{enumerate}
    \item Start from source vertex
    \item Add to queue and mark visited
    \item Dequeue, process, enqueue unvisited neighbors
    \item Repeat until queue empty
  \end{enumerate}

  \vspace{0.3cm}
  \textbf{Characteristics:}
  \begin{itemize}
    \item \textbf{Time}: $O(V + E)$
    \item \textbf{Space}: $O(V)$ for queue and visited set
    \item \textbf{Data Structure}: Queue (FIFO)
  \end{itemize}

  \vspace{0.3cm}
  \textbf{Key Property:}
  \begin{itemize}
    \item Finds \textbf{shortest path} in unweighted graphs
    \item Visits vertices in order of distance from source
  \end{itemize}
\end{frame}

\begin{frame}[fragile]{BFS Implementation}
\begin{lstlisting}[language=Python, basicstyle=\ttfamily\tiny]
from collections import deque

def bfs(graph, start):
    """BFS traversal from start node"""
    visited = set([start])
    queue = deque([start])
    result = []

    while queue:
        node = queue.popleft()
        result.append(node)

        for neighbor in graph[node]:
            if neighbor not in visited:
                visited.add(neighbor)
                queue.append(neighbor)

    return result

# Example
graph = {
    0: [1, 2],
    1: [0, 3, 4],
    2: [0, 4],
    3: [1],
    4: [1, 2]
}
print(bfs(graph, 0))  # [0, 1, 2, 3, 4]
\end{lstlisting}
\end{frame}

\begin{frame}{BFS Applications}
  \textbf{1. Shortest Path in Unweighted Graph:}
  \begin{itemize}
    \item Find minimum number of edges from source to target
    \item Track parent pointers to reconstruct path
  \end{itemize}

  \vspace{0.2cm}
  \textbf{2. Level-Order Traversal:}
  \begin{itemize}
    \item Process nodes by distance from source
    \item Used in tree level-order traversal
  \end{itemize}

  \vspace{0.2cm}
  \textbf{3. Connected Components:}
  \begin{itemize}
    \item Find all disconnected subgraphs
    \item Run BFS from each unvisited vertex
  \end{itemize}

  \vspace{0.2cm}
  \textbf{4. Bipartite Check:}
  \begin{itemize}
    \item Determine if graph is 2-colorable
    \item Alternate colors during BFS
  \end{itemize}

  \vspace{0.2cm}
  \textbf{5. Web Crawling:}
  \begin{itemize}
    \item Explore web pages level by level
    \item Limit crawl depth
  \end{itemize}
\end{frame}

\begin{frame}{Depth-First Search (DFS)}
  \textbf{Explore as Deep as Possible Using Stack/Recursion}

  \vspace{0.3cm}
  \textbf{Algorithm:}
  \begin{enumerate}
    \item Start from source vertex
    \item Mark visited, explore first unvisited neighbor
    \item Recursively DFS on neighbor
    \item Backtrack when no unvisited neighbors
  \end{enumerate}

  \vspace{0.3cm}
  \textbf{Characteristics:}
  \begin{itemize}
    \item \textbf{Time}: $O(V + E)$
    \item \textbf{Space}: $O(V)$ for recursion stack (or explicit stack)
    \item \textbf{Data Structure}: Stack (LIFO) or recursion
  \end{itemize}

  \vspace{0.3cm}
  \textbf{Key Property:}
  \begin{itemize}
    \item Explores entire branch before backtracking
    \item Can detect cycles (back edges)
  \end{itemize}
\end{frame}

\begin{frame}[fragile]{DFS Implementation}
\begin{lstlisting}[language=Python, basicstyle=\ttfamily\tiny]
def dfs_recursive(graph, node, visited=None):
    """DFS traversal (recursive)"""
    if visited is None:
        visited = set()

    visited.add(node)
    result = [node]

    for neighbor in graph[node]:
        if neighbor not in visited:
            result.extend(dfs_recursive(graph, neighbor, visited))

    return result

def dfs_iterative(graph, start):
    """DFS traversal (iterative)"""
    visited = set()
    stack = [start]
    result = []

    while stack:
        node = stack.pop()
        if node not in visited:
            visited.add(node)
            result.append(node)
            stack.extend(reversed(graph[node]))  # Maintain order

    return result
\end{lstlisting}
\end{frame}

\begin{frame}{DFS Applications}
  \textbf{1. Cycle Detection:}
  \begin{itemize}
    \item Find back edges (edge to ancestor)
    \item Detect if graph contains cycles
  \end{itemize}

  \vspace{0.2cm}
  \textbf{2. Path Finding:}
  \begin{itemize}
    \item Find any path between two nodes
    \item Find all paths (with backtracking)
  \end{itemize}

  \vspace{0.2cm}
  \textbf{3. Topological Sorting:}
  \begin{itemize}
    \item Order vertices in DAG
    \item Process finish times in reverse
  \end{itemize}

  \vspace{0.2cm}
  \textbf{4. Maze Solving:}
  \begin{itemize}
    \item Find path through grid
    \item Backtrack on dead ends
  \end{itemize}

  \vspace{0.2cm}
  \textbf{5. Strongly Connected Components:}
  \begin{itemize}
    \item Kosaraju's and Tarjan's algorithms
    \item Two-pass DFS
  \end{itemize}
\end{frame}

\begin{frame}{BFS vs DFS Comparison}
  \begin{center}
  \begin{tabular}{|l|c|c|}
    \hline
    \textbf{Feature} & \textbf{BFS} & \textbf{DFS} \\
    \hline
    Data Structure & Queue & Stack/Recursion \\
    \hline
    Path Found & Shortest & Any path \\
    \hline
    Memory (worst) & $O(V)$ (wide) & $O(h)$ (height) \\
    \hline
    Best For & Shortest path & Cycle detection \\
    \hline
    Tree Traversal & Level-order & Pre/In/Post-order \\
    \hline
    Completeness & Yes & Yes \\
    \hline
  \end{tabular}
  \end{center}

  \vspace{0.3cm}
  \textbf{When to Use BFS:}
  \begin{itemize}
    \item Find shortest path in unweighted graph
    \item Process nodes by distance
    \item Graph is very deep (avoid stack overflow)
  \end{itemize}

  \vspace{0.2cm}
  \textbf{When to Use DFS:}
  \begin{itemize}
    \item Detect cycles, find topological order
    \item Explore all paths, backtracking problems
    \item Graph is very wide (save memory)
  \end{itemize}
\end{frame}

\section{Shortest Paths: Dijkstra/Bellman-Ford}

\begin{frame}{Dijkstra's Algorithm}
  \textbf{Single-Source Shortest Path (Non-Negative Weights)}

  \vspace{0.3cm}
  \textbf{Algorithm:}
  \begin{enumerate}
    \item Initialize distances: source = 0, others = $\infty$
    \item Use min-heap to get vertex with minimum distance
    \item For each neighbor, relax edge if shorter path found
    \item Mark vertex as visited
    \item Repeat until all vertices processed
  \end{enumerate}

  \vspace{0.3cm}
  \textbf{Characteristics:}
  \begin{itemize}
    \item \textbf{Time}: $O((V + E) \log V)$ with min-heap
    \item \textbf{Space}: $O(V)$
    \item \textbf{Requirement}: Non-negative edge weights
  \end{itemize}

  \vspace{0.3cm}
  \textbf{Key Idea:}
  \begin{itemize}
    \item Greedy approach: always pick closest unvisited vertex
    \item Correctness requires non-negative weights
  \end{itemize}
\end{frame}

\begin{frame}[fragile]{Dijkstra Implementation}
\begin{lstlisting}[language=Python, basicstyle=\ttfamily\tiny]
import heapq

def dijkstra(graph, start):
    """Find shortest paths from start to all vertices"""
    dist = {node: float('inf') for node in graph}
    dist[start] = 0
    pq = [(0, start)]  # (distance, node)
    visited = set()

    while pq:
        d, node = heapq.heappop(pq)
        if node in visited:
            continue
        visited.add(node)

        for neighbor, weight in graph[node]:
            new_dist = d + weight
            if new_dist < dist[neighbor]:
                dist[neighbor] = new_dist
                heapq.heappush(pq, (new_dist, neighbor))

    return dist

# Example: graph[node] = [(neighbor, weight), ...]
graph = {
    'A': [('B', 4), ('C', 2)],
    'B': [('C', 1), ('D', 5)],
    'C': [('D', 8), ('E', 10)],
    'D': [('E', 2)],
    'E': []
}
print(dijkstra(graph, 'A'))  # {'A': 0, 'B': 4, 'C': 2, 'D': 9, 'E': 11}
\end{lstlisting}
\end{frame}

\begin{frame}{Bellman-Ford Algorithm}
  \textbf{Single-Source Shortest Path (Handles Negative Weights)}

  \vspace{0.3cm}
  \textbf{Algorithm:}
  \begin{enumerate}
    \item Initialize distances: source = 0, others = $\infty$
    \item Relax all edges $V-1$ times
    \item Check for negative cycles (one more iteration)
  \end{enumerate}

  \vspace{0.3cm}
  \textbf{Characteristics:}
  \begin{itemize}
    \item \textbf{Time}: $O(V \times E)$
    \item \textbf{Space}: $O(V)$
    \item \textbf{Advantage}: Handles negative weights, detects negative cycles
  \end{itemize}

  \vspace{0.3cm}
  \textbf{Why $V-1$ Iterations?}
  \begin{itemize}
    \item Longest simple path has $V-1$ edges
    \item Each iteration extends shortest path by one edge
    \item After $V-1$ iterations, all shortest paths found
  \end{itemize}
\end{frame}

\begin{frame}[fragile]{Bellman-Ford Implementation}
\begin{lstlisting}[language=Python, basicstyle=\ttfamily\tiny]
def bellman_ford(edges, n, start):
    """
    Find shortest paths, detect negative cycles
    edges: list of (u, v, weight)
    """
    dist = [float('inf')] * n
    dist[start] = 0

    # Relax edges V-1 times
    for _ in range(n - 1):
        for u, v, weight in edges:
            if dist[u] != float('inf') and dist[u] + weight < dist[v]:
                dist[v] = dist[u] + weight

    # Check for negative cycles
    for u, v, weight in edges:
        if dist[u] != float('inf') and dist[u] + weight < dist[v]:
            return None  # Negative cycle detected

    return dist

# Example
edges = [(0, 1, 4), (0, 2, 2), (1, 2, -3), (2, 3, 2), (3, 1, 1)]
n = 4
print(bellman_ford(edges, n, 0))
\end{lstlisting}
\end{frame}

\begin{frame}{Dijkstra vs Bellman-Ford}
  \begin{center}
  \begin{tabular}{|l|c|c|}
    \hline
    \textbf{Feature} & \textbf{Dijkstra} & \textbf{Bellman-Ford} \\
    \hline
    Time Complexity & $O(E \log V)$ & $O(V \times E)$ \\
    \hline
    Negative Weights & No & Yes \\
    \hline
    Negative Cycles & Cannot detect & Detects \\
    \hline
    Implementation & Min-heap & Nested loops \\
    \hline
    Best For & Fast, non-negative & Negative weights \\
    \hline
  \end{tabular}
  \end{center}

  \vspace{0.3cm}
  \textbf{Floyd-Warshall} (All-Pairs Shortest Paths):
  \begin{itemize}
    \item \textbf{Time}: $O(V^3)$
    \item Finds shortest paths between all pairs
    \item Handles negative weights
    \item Uses dynamic programming
  \end{itemize}

  \vspace{0.3cm}
  \textbf{Choosing Algorithm:}
  \begin{itemize}
    \item Non-negative weights $\rightarrow$ Dijkstra (faster)
    \item Negative weights $\rightarrow$ Bellman-Ford
    \item All pairs $\rightarrow$ Floyd-Warshall or run Dijkstra $V$ times
  \end{itemize}
\end{frame}

\section{Minimum Spanning Trees: Kruskal/Prim}

\begin{frame}{Minimum Spanning Tree (MST)}
  \textbf{Tree Connecting All Vertices with Minimum Total Weight}

  \vspace{0.3cm}
  \textbf{Properties:}
  \begin{itemize}
    \item Connects all $V$ vertices
    \item Has exactly $V-1$ edges
    \item No cycles (it's a tree)
    \item Minimum total edge weight
  \end{itemize}

  \vspace{0.3cm}
  \textbf{Applications:}
  \begin{itemize}
    \item Network design (minimize cable length)
    \item Approximation algorithms (TSP)
    \item Clustering algorithms
    \item Image segmentation
  \end{itemize}

  \vspace{0.3cm}
  \textbf{Two Main Algorithms:}
  \begin{itemize}
    \item \textbf{Kruskal's}: Sort edges, add if no cycle (edge-based)
    \item \textbf{Prim's}: Grow tree from vertex (vertex-based)
  \end{itemize}
\end{frame}

\begin{frame}{Kruskal's Algorithm}
  \textbf{Sort Edges, Add If Doesn't Create Cycle}

  \vspace{0.3cm}
  \textbf{Algorithm:}
  \begin{enumerate}
    \item Sort all edges by weight (ascending)
    \item Initialize Union-Find structure
    \item For each edge $(u, v)$:
    \begin{itemize}
      \item If $u$ and $v$ in different components, add edge
      \item Union the components
    \end{itemize}
    \item Stop when $V-1$ edges added
  \end{enumerate}

  \vspace{0.3cm}
  \textbf{Characteristics:}
  \begin{itemize}
    \item \textbf{Time}: $O(E \log E)$ (dominated by sorting)
    \item \textbf{Space}: $O(V)$ for Union-Find
    \item \textbf{Best for}: Sparse graphs
  \end{itemize}

  \vspace{0.3cm}
  \textbf{Key Data Structure:}
  \begin{itemize}
    \item Union-Find (Disjoint Set Union) for cycle detection
  \end{itemize}
\end{frame}

\begin{frame}[fragile]{Kruskal Implementation}
\begin{lstlisting}[language=Python, basicstyle=\ttfamily\tiny]
class UnionFind:
    def __init__(self, n):
        self.parent = list(range(n))
        self.rank = [0] * n

    def find(self, x):
        if self.parent[x] != x:
            self.parent[x] = self.find(self.parent[x])  # Path compression
        return self.parent[x]

    def union(self, x, y):
        px, py = self.find(x), self.find(y)
        if px == py:
            return False  # Already in same set
        if self.rank[px] < self.rank[py]:
            self.parent[px] = py
        elif self.rank[px] > self.rank[py]:
            self.parent[py] = px
        else:
            self.parent[py] = px
            self.rank[px] += 1
        return True

def kruskal(n, edges):
    """edges: list of (weight, u, v)"""
    edges.sort()  # Sort by weight
    uf = UnionFind(n)
    mst, total = [], 0

    for weight, u, v in edges:
        if uf.union(u, v):
            mst.append((u, v, weight))
            total += weight

    return total, mst
\end{lstlisting}
\end{frame}

\begin{frame}{Prim's Algorithm}
  \textbf{Grow Tree from Starting Vertex}

  \vspace{0.3cm}
  \textbf{Algorithm:}
  \begin{enumerate}
    \item Start with any vertex
    \item Add to MST
    \item Repeat until all vertices added:
    \begin{itemize}
      \item Find minimum weight edge from MST to non-MST vertex
      \item Add that edge and vertex to MST
    \end{itemize}
  \end{enumerate}

  \vspace{0.3cm}
  \textbf{Characteristics:}
  \begin{itemize}
    \item \textbf{Time}: $O(E \log V)$ with min-heap
    \item \textbf{Space}: $O(V)$
    \item \textbf{Best for}: Dense graphs
  \end{itemize}

  \vspace{0.3cm}
  \textbf{Key Data Structure:}
  \begin{itemize}
    \item Min-heap to efficiently find minimum edge
  \end{itemize}

  \vspace{0.3cm}
  \textbf{Similarity to Dijkstra:}
  \begin{itemize}
    \item Both use greedy approach with min-heap
    \item Prim: minimize edge weight, Dijkstra: minimize path distance
  \end{itemize}
\end{frame}

\begin{frame}[fragile]{Prim Implementation}
\begin{lstlisting}[language=Python, basicstyle=\ttfamily\tiny]
import heapq

def prim(graph, start):
    """
    graph: {node: [(neighbor, weight), ...]}
    """
    mst, total = [], 0
    visited = {start}
    edges = [(weight, start, neighbor)
             for neighbor, weight in graph[start]]
    heapq.heapify(edges)

    while edges:
        weight, u, v = heapq.heappop(edges)
        if v in visited:
            continue

        visited.add(v)
        mst.append((u, v, weight))
        total += weight

        for neighbor, w in graph[v]:
            if neighbor not in visited:
                heapq.heappush(edges, (w, v, neighbor))

    return total, mst

# Example
graph = {
    'A': [('B', 4), ('C', 2)],
    'B': [('A', 4), ('C', 1), ('D', 5)],
    'C': [('A', 2), ('B', 1), ('D', 8)],
    'D': [('B', 5), ('C', 8)]
}
print(prim(graph, 'A'))  # (7, [('A', 'C', 2), ('C', 'B', 1), ('B', 'D', 5)])
\end{lstlisting}
\end{frame}

\begin{frame}{Kruskal vs Prim}
  \begin{center}
  \begin{tabular}{|l|c|c|}
    \hline
    \textbf{Feature} & \textbf{Kruskal} & \textbf{Prim} \\
    \hline
    Time Complexity & $O(E \log E)$ & $O(E \log V)$ \\
    \hline
    Approach & Edge-based & Vertex-based \\
    \hline
    Data Structure & Union-Find & Min-heap \\
    \hline
    Best For & Sparse graphs & Dense graphs \\
    \hline
    Starting Point & N/A (all edges) & Any vertex \\
    \hline
    Works on Disconnected & Partial MST & No \\
    \hline
  \end{tabular}
  \end{center}

  \vspace{0.3cm}
  \textbf{Time Complexity Notes:}
  \begin{itemize}
    \item Kruskal: $O(E \log E) = O(E \log V)$ since $E \leq V^2$
    \item Prim with Fibonacci heap: $O(E + V \log V)$ (theoretical)
  \end{itemize}

  \vspace{0.3cm}
  \textbf{Practical Choice:}
  \begin{itemize}
    \item Sparse graph ($E \approx V$): Kruskal slightly better
    \item Dense graph ($E \approx V^2$): Prim slightly better
    \item Both give same MST weight (may differ in edges)
  \end{itemize}
\end{frame}

\section{Topological Sort and DAG DP}

\begin{frame}{Topological Sort}
  \textbf{Linear Ordering of Vertices in DAG}

  \vspace{0.3cm}
  \textbf{Definition:}
  \begin{itemize}
    \item For every directed edge $(u, v)$, $u$ comes before $v$ in ordering
    \item Only exists for Directed Acyclic Graphs (DAGs)
    \item Can be multiple valid orderings
  \end{itemize}

  \vspace{0.3cm}
  \textbf{Applications:}
  \begin{itemize}
    \item \textbf{Course scheduling}: Prerequisite dependencies
    \item \textbf{Build systems}: Compile dependencies (Makefile)
    \item \textbf{Task scheduling}: Task dependencies
    \item \textbf{Formula evaluation}: Dependency graphs
  \end{itemize}

  \vspace{0.3cm}
  \textbf{Two Main Algorithms:}
  \begin{itemize}
    \item \textbf{Kahn's Algorithm}: BFS-based, uses in-degrees
    \item \textbf{DFS-based}: Process vertices by finish time
  \end{itemize}
\end{frame}

\begin{frame}[fragile]{Topological Sort: Kahn's Algorithm}
\begin{lstlisting}[language=Python, basicstyle=\ttfamily\tiny]
from collections import deque

def topological_sort(graph, n):
    """
    Kahn's algorithm (BFS-based)
    graph: adjacency list
    Returns: topological order or None if cycle exists
    """
    # Calculate in-degrees
    in_degree = [0] * n
    for node in range(n):
        for neighbor in graph[node]:
            in_degree[neighbor] += 1

    # Start with vertices having in-degree 0
    queue = deque([i for i in range(n) if in_degree[i] == 0])
    result = []

    while queue:
        node = queue.popleft()
        result.append(node)

        for neighbor in graph[node]:
            in_degree[neighbor] -= 1
            if in_degree[neighbor] == 0:
                queue.append(neighbor)

    # If not all vertices processed, graph has cycle
    return result if len(result) == n else None
\end{lstlisting}
\end{frame}

\begin{frame}{DAG Dynamic Programming}
  \textbf{DP on Directed Acyclic Graphs}

  \vspace{0.3cm}
  \textbf{Key Idea:}
  \begin{itemize}
    \item Process vertices in topological order
    \item Each vertex computed after all dependencies
    \item No cycles $\rightarrow$ no circular dependencies
  \end{itemize}

  \vspace{0.3cm}
  \textbf{Common Problems:}
  \begin{itemize}
    \item \textbf{Longest/Shortest Path in DAG}: $O(V + E)$
    \item \textbf{Count paths}: Number of paths from source to sink
    \item \textbf{Critical Path Method}: Project scheduling
  \end{itemize}

  \vspace{0.3cm}
  \textbf{Template:}
  \begin{enumerate}
    \item Compute topological order
    \item Initialize DP array
    \item Process vertices in topological order
    \item Update DP based on edges
  \end{enumerate}
\end{frame}

\begin{frame}[fragile]{DAG DP: Longest Path}
\begin{lstlisting}[language=Python, basicstyle=\ttfamily\small]
def longest_path_dag(graph, n):
    """Find longest path in DAG"""
    topo_order = topological_sort(graph, n)
    if topo_order is None:
        return None  # Cycle exists

    dp = [0] * n

    for node in topo_order:
        for neighbor in graph[node]:
            dp[neighbor] = max(dp[neighbor], dp[node] + 1)

    return max(dp)

# Example: Course scheduling with prerequisites
# Find longest chain of courses
graph = {
    0: [1, 2],
    1: [3],
    2: [3],
    3: []
}
print(longest_path_dag(graph, 4))  # Output: 2
\end{lstlisting}
\end{frame}

\section{Strongly Connected Components}

\begin{frame}{Strongly Connected Components (SCC)}
  \textbf{Maximal Subgraphs Where Every Vertex Reaches Every Other}

  \vspace{0.3cm}
  \textbf{Definition:}
  \begin{itemize}
    \item In directed graph, SCC is maximal set of vertices
    \item For any $u, v$ in SCC, there exists path $u \rightarrow v$ and $v \rightarrow u$
    \item Condensation graph (SCC graph) is always a DAG
  \end{itemize}

  \vspace{0.3cm}
  \textbf{Applications:}
  \begin{itemize}
    \item \textbf{2-SAT}: Satisfiability of Boolean formulas
    \item \textbf{Reachability queries}: Which vertices can reach which
    \item \textbf{Deadlock detection}: Circular dependencies
    \item \textbf{Web page ranking}: Identify tightly connected clusters
  \end{itemize}

  \vspace{0.3cm}
  \textbf{Algorithms:}
  \begin{itemize}
    \item \textbf{Kosaraju's}: Two-pass DFS (simpler)
    \item \textbf{Tarjan's}: Single-pass DFS (more efficient)
  \end{itemize}
\end{frame}

\begin{frame}[fragile]{Kosaraju's Algorithm}
\textbf{Two-Pass DFS Approach}

\begin{lstlisting}[language=Python, basicstyle=\ttfamily\tiny]
def kosaraju_scc(graph, n):
    """Find strongly connected components"""
    # Step 1: DFS to get finish order
    visited = [False] * n
    finish_order = []

    def dfs1(node):
        visited[node] = True
        for neighbor in graph[node]:
            if not visited[neighbor]:
                dfs1(neighbor)
        finish_order.append(node)

    for i in range(n):
        if not visited[i]:
            dfs1(i)

    # Step 2: Transpose graph
    transpose = [[] for _ in range(n)]
    for u in range(n):
        for v in graph[u]:
            transpose[v].append(u)

    # Step 3: DFS on transpose in reverse finish order
    visited = [False] * n
    components = []

    def dfs2(node, component):
        visited[node] = True
        component.append(node)
        for neighbor in transpose[node]:
            if not visited[neighbor]:
                dfs2(neighbor, component)

    for node in reversed(finish_order):
        if not visited[node]:
            component = []
            dfs2(node, component)
            components.append(component)

    return components
\end{lstlisting}
\end{frame}

\begin{frame}{Kosaraju's Algorithm: How It Works}
  \textbf{Three Steps:}

  \vspace{0.3cm}
  \textbf{Step 1: First DFS}
  \begin{itemize}
    \item Run DFS on original graph
    \item Record finish times (when vertex fully explored)
    \item Vertices in same SCC finish close together
  \end{itemize}

  \vspace{0.3cm}
  \textbf{Step 2: Transpose Graph}
  \begin{itemize}
    \item Reverse all edge directions
    \item If $u \rightarrow v$ in $G$, then $v \rightarrow u$ in $G^T$
    \item SCCs remain the same
  \end{itemize}

  \vspace{0.3cm}
  \textbf{Step 3: Second DFS}
  \begin{itemize}
    \item Process vertices in reverse finish order
    \item Each DFS tree in $G^T$ is one SCC
  \end{itemize}

  \vspace{0.3cm}
  \textbf{Complexity:}
  \begin{itemize}
    \item Time: $O(V + E)$ (two DFS passes)
    \item Space: $O(V)$ for transpose graph
  \end{itemize}
\end{frame}

\section{Flow Algorithms Overview}

\begin{frame}{Maximum Flow Problem}
  \textbf{Find Maximum Flow from Source to Sink}

  \vspace{0.3cm}
  \textbf{Definition:}
  \begin{itemize}
    \item Given: Directed graph with edge capacities
    \item Find: Maximum amount of flow from source $s$ to sink $t$
    \item Constraints: Flow $\leq$ capacity, flow conservation
  \end{itemize}

  \vspace{0.3cm}
  \textbf{Key Concepts:}
  \begin{itemize}
    \item \textbf{Capacity}: Maximum flow on edge
    \item \textbf{Residual graph}: Remaining capacity after flow
    \item \textbf{Augmenting path}: Path with available capacity
    \item \textbf{Cut}: Partition of vertices into two sets
  \end{itemize}

  \vspace{0.3cm}
  \textbf{Max-Flow Min-Cut Theorem:}
  \begin{itemize}
    \item Maximum flow value = Minimum cut capacity
    \item Fundamental result in network flow theory
  \end{itemize}
\end{frame}

\begin{frame}{Ford-Fulkerson Method}
  \textbf{Find Augmenting Paths Until No More Exist}

  \vspace{0.3cm}
  \textbf{Algorithm:}
  \begin{enumerate}
    \item Initialize flow to 0
    \item While augmenting path exists:
    \begin{itemize}
      \item Find augmenting path (any path with capacity)
      \item Compute bottleneck capacity
      \item Add flow along path
      \item Update residual graph
    \end{itemize}
  \end{enumerate}

  \vspace{0.3cm}
  \textbf{Edmonds-Karp Algorithm:}
  \begin{itemize}
    \item Ford-Fulkerson with BFS for finding paths
    \item \textbf{Time}: $O(V \times E^2)$ (guaranteed polynomial)
    \item Always finds shortest augmenting path
  \end{itemize}

  \vspace{0.3cm}
  \textbf{Characteristics:}
  \begin{itemize}
    \item \textbf{Time}: $O(V \times E^2)$ for Edmonds-Karp
    \item \textbf{Space}: $O(V^2)$ for flow/capacity matrices
  \end{itemize}
\end{frame}

\begin{frame}[fragile]{Edmonds-Karp Implementation}
\begin{lstlisting}[language=Python, basicstyle=\ttfamily\tiny]
from collections import deque

def max_flow(capacity, source, sink):
    """Edmonds-Karp algorithm for maximum flow"""
    n = len(capacity)
    flow = [[0] * n for _ in range(n)]

    def bfs():
        """Find augmenting path using BFS"""
        parent = [-1] * n
        visited = [False] * n
        visited[source] = True
        queue = deque([(source, float('inf'))])

        while queue:
            u, min_cap = queue.popleft()

            for v in range(n):
                if not visited[v] and capacity[u][v] - flow[u][v] > 0:
                    visited[v] = True
                    parent[v] = u
                    new_cap = min(min_cap, capacity[u][v] - flow[u][v])

                    if v == sink:
                        return parent, new_cap
                    queue.append((v, new_cap))

        return None, 0

    total_flow = 0
    while True:
        parent, path_flow = bfs()
        if path_flow == 0:
            break

        total_flow += path_flow
        v = sink
        while v != source:
            u = parent[v]
            flow[u][v] += path_flow
            flow[v][u] -= path_flow  # Reverse flow
            v = u

    return total_flow
\end{lstlisting}
\end{frame}

\begin{frame}{Flow Applications}
  \textbf{1. Maximum Bipartite Matching:}
  \begin{itemize}
    \item Model as flow problem
    \item Add source to left vertices, sink from right vertices
    \item Maximum flow = maximum matching
  \end{itemize}

  \vspace{0.2cm}
  \textbf{2. Minimum Cut:}
  \begin{itemize}
    \item Find minimum capacity cut separating $s$ and $t$
    \item Max flow = min cut (by theorem)
    \item Applications: Network reliability, image segmentation
  \end{itemize}

  \vspace{0.2cm}
  \textbf{3. Network Routing:}
  \begin{itemize}
    \item Optimize data flow through network
    \item Consider bandwidth constraints
  \end{itemize}

  \vspace{0.2cm}
  \textbf{4. Assignment Problems:}
  \begin{itemize}
    \item Match workers to tasks optimally
    \item Constraint satisfaction
  \end{itemize}

  \vspace{0.2cm}
  \textbf{5. Circulation with Demands:}
  \begin{itemize}
    \item Flow with lower and upper bounds
    \item Supply and demand at vertices
  \end{itemize}
\end{frame}

\section{Summary}

\begin{frame}{Key Takeaways}
  \textbf{Traversal Algorithms:}
  \begin{itemize}
    \item \textbf{BFS}: Shortest path, level-order, $O(V + E)$
    \item \textbf{DFS}: Cycle detection, topological sort, $O(V + E)$
  \end{itemize}

  \vspace{0.3cm}
  \textbf{Shortest Path Algorithms:}
  \begin{itemize}
    \item \textbf{Dijkstra}: Non-negative weights, $O(E \log V)$
    \item \textbf{Bellman-Ford}: Negative weights, detects cycles, $O(VE)$
  \end{itemize}

  \vspace{0.3cm}
  \textbf{Minimum Spanning Tree:}
  \begin{itemize}
    \item \textbf{Kruskal}: Sort edges, Union-Find, $O(E \log E)$
    \item \textbf{Prim}: Grow tree, min-heap, $O(E \log V)$
  \end{itemize}

  \vspace{0.3cm}
  \textbf{Advanced Topics:}
  \begin{itemize}
    \item \textbf{Topological Sort}: Order DAG vertices, $O(V + E)$
    \item \textbf{SCC}: Kosaraju's/Tarjan's, $O(V + E)$
    \item \textbf{Max Flow}: Edmonds-Karp, $O(VE^2)$
  \end{itemize}
\end{frame}

\begin{frame}{Complexity Summary}
  \begin{center}
  \begin{tabular}{|l|c|c|}
    \hline
    \textbf{Algorithm} & \textbf{Time} & \textbf{Space} \\
    \hline
    BFS & $O(V + E)$ & $O(V)$ \\
    \hline
    DFS & $O(V + E)$ & $O(V)$ \\
    \hline
    Dijkstra & $O((V+E) \log V)$ & $O(V)$ \\
    \hline
    Bellman-Ford & $O(VE)$ & $O(V)$ \\
    \hline
    Kruskal & $O(E \log E)$ & $O(V)$ \\
    \hline
    Prim & $O(E \log V)$ & $O(V)$ \\
    \hline
    Topological Sort & $O(V + E)$ & $O(V)$ \\
    \hline
    Kosaraju SCC & $O(V + E)$ & $O(V)$ \\
    \hline
    Edmonds-Karp & $O(VE^2)$ & $O(V^2)$ \\
    \hline
  \end{tabular}
  \end{center}
\end{frame}

\begin{frame}{Practice Problems}
  \textbf{BFS/DFS:}
  \begin{itemize}
    \item Number of islands (LeetCode 200)
    \item Word ladder (LeetCode 127)
    \item Course schedule (LeetCode 207, 210)
  \end{itemize}

  \vspace{0.2cm}
  \textbf{Shortest Paths:}
  \begin{itemize}
    \item Network delay time (LeetCode 743)
    \item Cheapest flights (LeetCode 787)
  \end{itemize}

  \vspace{0.2cm}
  \textbf{MST:}
  \begin{itemize}
    \item Min cost to connect all points (LeetCode 1584)
  \end{itemize}

  \vspace{0.2cm}
  \textbf{Advanced:}
  \begin{itemize}
    \item Critical connections (LeetCode 1192)
    \item Alien dictionary (LeetCode 269)
  \end{itemize}
\end{frame}

\begin{frame}{Resources}
  \textbf{Books:}
  \begin{itemize}
    \item "Introduction to Algorithms" (CLRS) - Chapters 22-26
    \item "Algorithm Design" (Kleinberg \& Tardos)
  \end{itemize}

  \vspace{0.3cm}
  \textbf{Online:}
  \begin{itemize}
    \item VisuAlgo - Graph algorithm visualizations
    \item LeetCode - Graph problems
    \item Codeforces - Graph theory tutorials
  \end{itemize}

  \vspace{0.3cm}
  \textbf{Advanced Topics:}
  \begin{itemize}
    \item A* search algorithm
    \item Network simplex for min-cost flow
    \item Hopcroft-Karp for bipartite matching
    \item Tarjan's algorithm for SCC
  \end{itemize}
\end{frame}

\end{document}
